\documentclass[12pt]{article}

\usepackage[top=2.2in, left=0.75in,right=0.75in,headsep=1mm,footskip=.5in]{geometry}
\input{document_includes}

\pagestyle{fancy}
\let\headrule\empty
\let\footrule\empty
\rhead{\hrule}
\chead{\hrule}
\lhead{\hrule}
\lfoot{CSC\,263\,H1, Fall 2025}
\cfoot{Page \thepage\space of \pageref{LastPage}}
\rfoot{Quiz 3}

\begin{document}

\begin{center}\large
UNIVERSITY OF TORONTO\\
Faculty of Arts and Science\\[1.75ex]
Department of Computer Science\\[1.75ex]
\textbf{Quiz 3} \\
\textbf{{CSC\,263\,H1}}
\end{center}

\infobox[\textbf{First Name:}]
\infobox[\textbf{Last Name:}]

\noindent \textbf{Student \# (9 or 10 digits):} \begin{table}[H]
\begin{tabular}{|c|c|c|c|c|c|c|c|c|c|}
\hline
\hspace{1.2em}\; & \hspace{1.2em}\; & \hspace{1.2em}\; & \hspace{1.2em}\; & \hspace{1.2em}\; & \hspace{1.2em}\; & \hspace{1.2em}\; & \hspace{1.2em}\; & \hspace{1.2em}\; & \hspace{1.2em}\; \\
\hspace{1.2em}\; & \hspace{1.2em}\; & \hspace{1.2em}\; & \hspace{1.2em}\; & \hspace{1.2em}\; & \hspace{1.2em}\; & \hspace{1.2em}\; & \hspace{1.2em}\; & \hspace{1.2em}\; & \hspace{1.2em}\;\\\hline
\end{tabular}%
\end{table}
\noindent \textbf{UTORid (e.g. \texttt{gummalur}):} \begin{table}[H]
\begin{tabular}{|c|c|c|c|c|c|c|c|}
\hline
\hspace{1.2em}\; & \hspace{1.2em}\; & \hspace{1.2em}\; & \hspace{1.2em}\; & \hspace{1.2em}\; & \hspace{1.2em}\; & \hspace{1.2em}\; & \hspace{1.2em}\; \\
\hspace{1.2em}\; & \hspace{1.2em}\; & \hspace{1.2em}\; & \hspace{1.2em}\; & \hspace{1.2em}\; & \hspace{1.2em}\; & \hspace{1.2em}\; & \hspace{1.2em}\;\\\hline
\end{tabular}%
\end{table}
\noindent \textbf{Duration:}\dotfilltext {50 minutes}\\
\textbf{Instructor:}\dotfilltext Chandra Gummaluru\\
\textbf{Weight:}\dotfilltext 10\% \\
%\textbf{Aids Allowed:}
%\begin{enumerate}
%    \item[(1)] Non-Programmable Calculator 
%\end{enumerate}
\\~\\
\begin{center}
\textit{Do NOT turn over the page until you are told to do so.}
\end{center}

\clearpage

\setlength\parindent{0pt}
\subsection*{Question 3.1: Unsorted Dictionaries}
Consider a \texttt{Dictionary} of size 5, which uses the hash function defined below:
\[
h(k) = k \bmod 5
\]
Suppose the following operations are done in order. For each operation, indicate the number of indices that will be visited before the operation is done. \\

\texttt{insert(12, val1)}: \numbox[2cm] \\
\texttt{insert(27, val2)}: \numbox[2cm] \\
\texttt{insert(33, val3)}: \numbox[2cm] \\ 
\texttt{insert(44, val4)}: \numbox[2cm] \\ 
\texttt{insert(15, val5)}: \numbox[2cm] \\
\texttt{delete(27)}:\hspace{13mm} \numbox[2cm] \\
\texttt{delete(44)}:\hspace{13mm} \numbox[2cm]
\\~\\
\textit{Note 1:} During insertion, probing continues until an empty slot is found. If the search returns to the original starting index, the table is considered full and the operation stops.
\\[2ex]
\textit{Note 2:} When an element is deleted, it is replaced with a special marker called a \texttt{pass} flag. During probing (for insertion, deletion, or search), continue while encountering \texttt{pass} flags, until the target element is found, an empty slot is reached, or the probe returns to the original starting index.

\newpage

\subsection*{Question 3.2: Tries}
\points[4] Suppose we use a trie to store DNA sequences (strings composed of \texttt{A}, \texttt{T}, \texttt{C}, and \texttt{G}) in order to identify genetic similarities.
\begin{enumerate}
\item[(a)] \points[3] Draw the trie after inserting the following DNA sequences in the specified order: \texttt{ATGCA}, \texttt{ATGCC}, \texttt{ATGTA}, \texttt{ATCGA}, \texttt{CGTAA}
\vspace{12cm}
\item[(b)] \points[1] Using the constructed trie, determine and illustrate the longest prefix match for the query sequence \texttt{ATGTCG}. Highlight the path on your tree from (a).
\end{enumerate}

\newpage 


\subsection*{Question 3.3: Queues}
\points[4] Consider a \texttt{Queue}, \texttt{queue} (built using a binary max-heap) with 5 elements with priorities; \texttt{10}, \texttt{25}, \texttt{30}, \texttt{40}, and \texttt{70}.
\begin{enumerate}
\item[(a)]\points[2] Draw a tree for the underlying binary max-heap (assuming the elements were inserted in the order shown above).
\vspace{6cm}
\item[(b)] \points[2] Suppose we insert an element with priority \texttt{45} (i.e., we call \texttt{queue.insert(obj, 45)}), and then pop an element (i.e., we call \texttt{queue.pop()}). Draw a tree for the underlying binary max-heap after both operations.
\end{enumerate}

\clearpage

\subsection*{Question 3.4: Sorted Dictionaries}
\points[4] Consider a \texttt{Dictionary}, \texttt{dict}, represented by the balanced binary search tree shown below:
\begin{center}
\begin{tikzpicture}[scale=0.8, transform shape]
\node[defaultnode] (n1) at (0,0)  { \texttt{28} };
\node[defaultnode] (n2) at (-3,-1.5)  { \texttt{21} };
\node[defaultnode] (n3) at (3,-1.5)  { \texttt{41} };
\node[defaultnode] (n4) at (1.5,-3)  { \texttt{35} };
\node[defaultnode] (n5) at (4.5,-3)  { \texttt{47} };
\node[defaultnode] (n6) at (-4.5,-3) { \texttt{18} };
\node[defaultnode] (n7) at (-1.5,-3) { \texttt{25} };
\node[defaultnode] (n8) at (0.5,-4.5) { \texttt{31} };

\draw[->] (n1) -- (n2);
\draw[->] (n1) -- (n3);
\draw[->] (n2) -- (n6);
\draw[->] (n2) -- (n7);
\draw[->] (n3) -- (n4);
\draw[->] (n3) -- (n5);
\draw[->] (n4) -- (n8);
\end{tikzpicture}
\end{center}

Suppose a new node with key \texttt{33} is inserted into the dictionary, i.e., we call \texttt{dict.insert(33)}.
\begin{itemize}
\item[(a)]\points[1] Indicate where the node would be attached in the tree (before re-balancing) by drawing the new node in the figure above.
\item[(b)]\points[1] After (a), which node in the tree is an anchor node?
\numbox[2cm] \checkbox[None]
\item[(c)]\points[1] Which rotation is needed to re-balance the tree? \\[2mm]
\checkbox[Left] \checkbox[Right] \checkbox[Left-Right] \checkbox[Right-Left] \checkbox[None]
\item[(d)] \points[1] Draw the resulting tree (after re-balancing).
\end{itemize}

\clearpage

\begin{center}
\textit{(Rough work — this page will not be marked)}
\end{center}
\clearpage

\end{document} 